% Results scaffold (Phase 2): five subsections; opens with one-sentence
% structural map. Target ~1700 w total. Each subsection has named-paragraph
% stubs, equation placeholders where required, and figure placeholders for
% the three planned main-body figures.
%
% Each subsection's content is project-status-dependent (more campaigns to
% come). Phase 2 fixes structure; content moves into prose during the
% per-section drafting session, drawing on whatever campaign data exists at
% that moment.

\section{Results}%
\label{sec:Results}

% One-sentence structural map.
% TODO: open with the §4-internal roadmap per genre_conventions_theory.md
% §(g) Pattern 5: "§4.1 validates TIDAL against the Boccaletti formula;
% §4.2–4.3 report sector-level constraints from the T1 and T2 channel
% surveys; §4.4 presents the propagating-PGT structural finding;
% §4.5 reports the trace-channel ghost diagnosis." (One sentence.)

\subsection{Einstein--Maxwell baseline}%
\label{sec:Validation}
% Target ~250 w. Leads with the formula; figure 1 placeholder; the actual
% percentage agreement will be filled during drafting from the final
% Boccaletti validation campaign.

\paragraph*{Validation setup and result}%
\label{par:ValidationSetup}
% TODO: state the validation point ($\Bzero$, $D$, $t_\text{end}$, grid
% resolution); cross-ref \cref{eq:GertsenshteinFormula}; report the
% measured agreement quantitatively (TBD: percentage from final HPC run);
% frame as a physics result, not a software test.

% Figure 1 — Boccaletti validation.
\begin{figure}[t]
  \centering
  % TODO: \includegraphics{fig1_boccaletti_validation}
  % TODO caption: "Simulated vs analytic conversion probability $\Pconv(t)$
  % for Einstein–Maxwell on a flat Minkowski background with uniform
  % $\Bzero = [\text{value}]$ and propagation distance $D = [\text{value}]$.
  % Upper panel: simulated (solid) and analytic (dashed); lower panel:
  % residual. The maximum relative deviation across the validation grid is
  % $[\text{value}]\%$, demonstrating the TIDAL pipeline reproduces the
  % Boccaletti kernel."
  \caption{TODO: Boccaletti validation; sim vs analytic with residual.}%
  \label{fig:BoccalettiValidation}
\end{figure}

\subsection{T1 sector: dark-photon--plasma channels}%
\label{sec:T1null}
% Target ~350 w.

\paragraph*{Constraint statement}%
\label{par:T1Constraint}
% TODO: Dandoy precision-bound framing — "we constrain the T1 sector
% conversion amplification to $A < [\text{value}]$ across the
% $(m_A, \varepsilon)$ parameter plane." Cite \cite{caputo2021darkphoton}
% for the dark-photon parameter context.

\paragraph*{Physical mechanism}%
\label{par:T1Mechanism}
% TODO: one sentence on the physical mechanism — what causes the no-amplification
% finding ("the [coupling term] vanishes at linear order because…").
% Cross-ref \cref{fig:StageAHeatmap} (App D).

\subsection{T2 sector: Einstein--Cartan parameter sweep}%
\label{sec:T2null}
% Target ~350 w.

\paragraph*{Constraint statement}%
\label{par:T2Constraint}
% TODO: same framing template — "we constrain the T2 sector conversion
% amplification to $A < [\text{value}]$ across the [parameter region]."
% State the number of runs (276 in the d21_barker_sup campaign).

\paragraph*{Physical mechanism}%
\label{par:T2Mechanism}
% TODO: one sentence on the physical mechanism. Cross-ref
% \cref{fig:StageBSweep} (App D).

\subsection{Propagating PGT: $h_\times \leftrightarrow \mathrm{ax}$ torsion-independence}%
\label{sec:PropagatingPGT}
% Target ~400 w. Competing-effects framing per genre_conventions_theory.md
% §(g) Pattern 3.

\paragraph*{Structural observation}%
\label{par:HxAxObservation}
% TODO: state the result as a structural observation visible in the
% symbolic output, not a derivation: "The hx and ax torsion sectors
% enter the kinetic matrix \cref{eq:KineticMatrix} with identical
% coefficients (\cref{eq:HxAxIdentity}, \cref{app:OutputsPropagating})."
% Frame consequence: "Consequently, the conversion probability is
% torsion-independent in this sector — agreeing with the Einstein–Maxwell
% prediction across the full stability window $|\delta_1| < 0.005$."

\paragraph*{Numerical confirmation}%
\label{par:HxAxNumerics}
% TODO: state the numerical finding ($A = 1$ to $[\text{tolerance}]$);
% cross-ref figure 2.

% Figure 2 — Propagating-torsion parameter map (the headline result figure).
\begin{figure}[t]
  \centering
  % TODO: \includegraphics{fig2_propagating_torsion_map}
  % TODO caption: "Amplification factor $A = \Pconv^\text{torsion}/\Pconv^\text{GR}$
  % vs torsion stability parameter $\delta_1$ for the propagating-PGT model
  % at $\Bzero = [\text{value}]$ and $t_\text{end} = [\text{value}]$. $A = 1$
  % across the full stability window $|\delta_1| < 0.005$ demonstrates the
  % $h_\times \leftrightarrow \mathrm{ax}$ torsion-independence
  % (\cref{eq:HxAxIdentity})."
  \caption{TODO: propagating-torsion parameter map; $A$ vs $\delta_1$.}%
  \label{fig:PropagatingTorsionMap}
\end{figure}

\subsection{Trace-channel ghost diagnosis}%
\label{sec:GhostDiagnosis}
% Target ~200 w (downscoped from 350 w because there is no dedicated
% appendix-level investigation behind this finding). Three-step ghost
% reporting per genre_conventions_theory.md §(g) Pattern 4.
%
% Decision deferred to drafting session: include \cref{fig:TraceGhostGrowth}
% if data quality supports a clean exponential-growth signal; otherwise
% drop fig 3 and reduce the subsection to ~150 w.

\paragraph*{No-ghost condition}%
\label{par:NoGhostCondition}
% TODO: state the closed inequality on the relevant coupling combination;
% cite \cite{sezgin1980ghostfree} and \cite{lin2019pgt} for the
% no-ghost-condition provenance.
\begin{align}
  % TODO: no-ghost inequality
  \label{eq:NoGhostCondition}
\end{align}

\paragraph*{Single-point diagnosis}%
\label{par:GhostDiagnosis}
% TODO: report whether the inequality is satisfied at the studied
% parameter point; cross-ref \cref{app:OutputsTrace} for the kinetic
% matrix at that point. State the consequence with caveat: "We diagnose
% the trace channel as ghost-unstable at this single parameter point;
% this is not a systematic claim about the trace sector across PGT."

% Figure 3 — Trace-channel ghost growth (CONDITIONAL — drafting session
% decides whether to include based on data quality).
\begin{figure}[t]
  \centering
  % TODO conditional: \includegraphics{fig3_trace_ghost_growth}
  % TODO caption: "Time evolution of $\log|h_\text{trace}(t)|$ at the
  % parameter point $[\text{values}]$. The exponential growth identifies
  % the trace channel as ghost-unstable at this point."
  \caption{TODO: trace-channel ghost growth (conditional).}%
  \label{fig:TraceGhostGrowth}
\end{figure}
